% Public CV: phone, grades, test scores, referees and hidden results removed from the private CV.
% Rebuild: npm run cv:pdf
% Document class and font size
\documentclass[a4paper,11pt]{extarticle}

% Packages
\usepackage[utf8]{inputenc}
\usepackage[T1]{fontenc}
\usepackage{mathptmx}
\usepackage[a4paper, margin=0.55in]{geometry}
\usepackage{titlesec}
\usepackage{needspace}
\usepackage{enumitem}
\usepackage{xcolor}
\usepackage{hyperref}

% Formatting
\definecolor{linkblue}{RGB}{60,100,135}
\setlist{noitemsep, topsep=1pt, leftmargin=*}
\titleformat{\section}{\needspace{5\baselineskip}\large\bfseries}{}{0em}{}[\vspace{-2pt}\titlerule]
\titlespacing*{\section}{0pt}{0.7\baselineskip}{0.4\baselineskip}
\hypersetup{colorlinks=true, urlcolor=linkblue, linkcolor=linkblue}
\pagestyle{empty}
\setlength{\parindent}{0pt}

% CV entry macros
\newlength{\datewidth}
\setlength{\datewidth}{2.1cm}

% Push #1 to the right margin, moving it to its own line if it would collide
\newcommand{\rightlabel}[1]{%
    \unskip\nobreak\hfil\penalty50\hskip0.8em\hbox{}\nobreak\hfill\mbox{#1}}

% \cventry{date}{title}{right-aligned label}
\newcommand{\cventry}[3]{%
    \par\addvspace{5pt}\begingroup\raggedright\hyphenpenalty=10000\exhyphenpenalty=10000
    \noindent\hangindent=\datewidth\hangafter=1
    \makebox[\datewidth][l]{\textit{#1}}\textbf{\MakeUppercase{#2}}\rightlabel{\textbf{#3}}\par\endgroup}

% Indented body text under an entry
\newenvironment{cvbody}
    {\list{}{\leftmargin=\datewidth\topsep=1pt\partopsep=0pt\parsep=0pt\itemsep=0pt}\item[]}
    {\endlist}

% \referee{name}{role}{email}{position}
\newcommand{\referee}[4]{%
    \par\addvspace{6pt}\noindent\textbf{#1} \textit{(#2)}\hfill\href{mailto:#3}{\textit{#3}}\\
    \textit{#4}\par}

\begin{document}

% Header
\begin{center}
    {\LARGE \textbf{BISHESH KHANAL}}\\[3pt]
    \textit{\href{mailto:Bisheshkhanal036@gmail.com}{Bisheshkhanal036@gmail.com} \; | \; \href{https://bisheshkhanal036-afk.github.io}{bisheshkhanal036-afk.github.io} \; | \; \href{https://www.linkedin.com/in/bisheshkhanal/}{LinkedIn}}
\end{center}

\vspace{2pt}
Civil Engineering graduate with first-author research in flood inundation modelling and precipitation extremes. Experienced in HEC-RAS 2D, HEC-HMS, GIS, and Python-based machine learning for hydrological prediction.

% Education
\section*{Education}
\cventry{2022--2026}{Tribhuvan University (TU)}{KATHMANDU, NP}
\begin{cvbody}
    Bachelor of Civil Engineering, IOE Thapathali Campus, TU
\end{cvbody}

\cventry{2020--2022}{Prasadi Academy}{LALITPUR, NP}
\begin{cvbody}
    High school education
\end{cvbody}

% Research Interests
\section*{Research Interests}
Flood risk modelling and hydraulic simulation, machine learning surrogates for real-time flood early warning, precipitation and streamflow extremes, and climate-risk assessment in data-sparse mountainous regions.

% Research and Manuscripts
\section*{Projects \& Publications}

\cventry{2025--2026}{Flood Inundation Mapping and Evacuation Routing: An Approach for Flood Risk Management}{FINAL YEAR PROJECT}
\begin{cvbody}
    \textit{(under review, Environmental Development, Elsevier; first author)}
    \begin{itemize}
        \item Developed a hydrologic and hydraulic model and built 2D flood simulation models to analyze spatial flood extent and flow dynamics
        \item Analyzed the critical infrastructure within the flood hazard zones
        \item Proposed evacuation strategies based on flood depth and hazard zones
    \end{itemize}
\end{cvbody}

\cventry{Aug 2026}{Rainfall Variability and Extreme Event Classification Using Machine Learning in the Narayani Basin, Nepal}{}
\begin{cvbody}
    \textit{(under review, Journal of Water and Climate Change, first author)}
\end{cvbody}

\cventry{June 2026}{Comparing Process-Based and Machine Learning Models for Streamflow Prediction in the Kaligandaki River Basin, Nepal}{EARTHARXIV}
\begin{cvbody}
    \textit{DOI of preprint:} \href{https://doi.org/10.31223/X55N3N}{10.31223/X55N3N}
    \begin{itemize}
        \item Developed and evaluated machine-learning streamflow prediction models in Python
        \item Applied feature engineering, hyperparameter tuning, and SHAP-based explainability techniques
        \item Compared process-based and data-driven approaches for hydrological prediction and forecasting
    \end{itemize}
\end{cvbody}

\cventry{2025}{Flood Hazard Mapping and Structural Design Recommendation}{8TH SEMESTER PROJECT}
\begin{cvbody}
    \begin{itemize}
        \item Conducted floodplain surveys along the Hanumante and Manohara river systems to support flood hazard mapping and hydraulic analysis
        \item Applied the surveyed data to generate terrain for detailed analysis
        \item Generated inundation maps and applied HEC-RAS and HEC-HMS for hydraulic and hydrologic simulation
        \item Recommended both engineering and non-engineering measures for reducing flood vulnerability
    \end{itemize}
\end{cvbody}

% Technical Skills
\section*{Technical Skills}
\textbf{Programming:} \textit{Python, LaTeX}\\[3pt]
\textbf{Libraries \& Frameworks:} \textit{NumPy, Pandas, Scikit-learn, PyTorch, XGBoost, SHAP}\\[3pt]
\textbf{Modelling \& GIS Tools:} \textit{HEC-RAS 2D, HEC-HMS, SWAT, ArcGIS Pro, QGIS, ArcMap}\\[3pt]
\textbf{ML \& DL Models:} \textit{Random Forest, XGBoost, SVM, LSTM, GRU}\\[3pt]
\textbf{Domains:} \textit{Flood \& Hydrological Modelling, Precipitation Extremes, Water Resources, Remote Sensing \& GIS}\\[3pt]
\textbf{Languages:} \textit{Nepali (native), English, Hindi}

% Courses
\section*{Certifications \& Professional Courses}
\begin{itemize}[label=\textbullet]
    \item \textbf{Supervised Machine Learning: Regression and Classification} --- DeepLearning.AI \& Stanford \rightlabel{\href{https://coursera.org/share/e5669465671a34e5445bc2617c6c8c2a}{Credential}}
    \item \textbf{Advanced Learning Algorithms} --- DeepLearning.AI \& Stanford
    \item \textbf{Neural Networks and Deep Learning} --- DeepLearning.AI \rightlabel{\href{https://coursera.org/share/072eb4c562e57e2f11b669c9afc13798}{Credential}}
\end{itemize}

% Conferences & Workshops
\section*{Conferences \& Workshops}
\cventry{Dec 2026}{Co-author Abstract, H091 (Integrative Approaches to Urban Flooding)}{AGU26}
\begin{cvbody}
    ``A Deep Learning Surrogate of a HEC-RAS 2D Model for Urban Flood Forecasting in the Kathmandu Valley, Nepal'' \\
    Trained a GRU surrogate of a calibrated HEC-RAS 2D model on 244 synthetic flood events; achieved an overall R\textsuperscript{2} of 0.947 and mean NSE above 0.88 at a 3-hour lead time. \textit{(abstract accepted, AGU Fall Meeting 2026, San Francisco)}
\end{cvbody}

\cventry{Aug 2026}{Hands-on Training Workshop on Measuring Urban Water Security}{LALITPUR, NP}
\begin{cvbody}
    \textit{ABCD Centre, Asian Institute of Technology (AIT) --- with IOE, TU, and CREEW, Nepal}
\end{cvbody}

\cventry{June 2026}{Presenter, 4\textsuperscript{th} Thapathali Graduate Conference}{THAPATHALI, NP}
\begin{cvbody}
    ``Flood Inundation Mapping and Evacuation Routing: An Approach for Flood Risk Management''
\end{cvbody}

% Honors and Awards
\section*{Honors and Awards}
\cventry{2026}{Best Final Year Project}{IOE THAPATHALI}
\begin{cvbody}
    Recognised as one of the best final year projects by the department project committee.
\end{cvbody}

\par\addvspace{5pt}\noindent\makebox[\datewidth][l]{\textit{2022--2026}}\textbf{SEMESTER EXCELLENCE UNDERGRAD AWARD} \textit{(received 6 times)}\hfill\textbf{IOE, TU}\par

\vspace{8pt}
\textit{References available on request.}

\end{document}
